\SourceLeaf{216}$cbe = cb'b$. Also wird $be$ die 2 in $e'$ schneiden. Macht man nun $b'g = be'$, so bilden $cg$ und $cd'$ mit $cb'$ gleiche Winkel. Das ist widersinnig.

5. Wenn 1 zu 2 und 1 zu 3 parallel ist, so ist auch 2 zu 3 parallel.

6. Was korrespondierende Punkte auf zwei Parallelen sind.

7. Gleicher Abstand der korrespondierenden Punkte.

8. Der Punkt auf einer dritten Parallelen korrespondiert mit korrespondierenden Punkten auf den beiden ersten.

9. \textls{Trope} ist die L[inie, die von korrespondierenden Punkten gebildet wird, wenn man alle Parallelen zu einer Geraden betrachtet.]

\Needspace{19\baselineskip}
\bigskip
\begin{center}
\rule{2.5cm}{0.4pt}

\textsc{Bemerkung.}
\end{center}

In dem Brief an \textsc{Schumacher} vom 17. Mai 1831 (S. 213 dieses Bandes) sagt \textsc{Gauss}, dass er von seinen Überlegungen zur Theorie der Parallellinien, die schon etwa 40 Jahre alt seien, zuvor nie etwas aufgeschrieben habe, dass er jedoch vor einigen Wochen begonnen habe, einiges aufzuschreiben. Auch in dem Brief an \textsc{Bolyai} vom 6. März 1832 erwähnt \textsc{Gauss} solche Aufzeichnungen. Man wird daher wohl nicht fehlgehen, wenn man annimmt, dass die vorstehenden undatierten Notizen auf einzelnen Zetteln aus dem Jahre 1831 stammen.

Die Figuren in der Notiz [1, sowie die erste Figur in der Notiz [3] sind dem Text hinzugefügt worden

\hfill \textsc{Stäckel}.

